\section{The instrument}
\label{sec:instrument}

Every result in this paper was found in a drawing tool before it was
proved, and measured in the tool after. \Moire\ is a web application that
draws superpositions of families the way this paper describes them. Each
layer is a count $\cnt_i$ and a stroke profile; at every pixel the count is
evaluated --- by a formula for the exact and winding rungs, by a certified
search for the fold rung --- and the picture on the torus is looked up
there and then. No layer is ever drawn into an intermediate image and
resampled, so a superposition has no resolution of its own: zooming asks
the same counts at new points, and cannot invent a beat that the geometry
does not have or destroy one it does. Fourteen families ship, sorted by
rung: on the exact rung, parallel lines, concentric circles, squares,
triangles and hexagons, the Archimedean and the logarithmic spiral, the
wave, the parabola and the hyperbola, and three lattices (square,
triangular, hexagonal) with their tilings; on the winding rung, the rays;
and on the fold rung, the walking versions of the concentric shapes. A
field language adds any function of position to a count.

\begin{figure}[t]
  \centering
  
\includegraphics[width=0.32\linewidth]{instrument-render.png}\hfill
  
\includegraphics[width=0.32\linewidth]{instrument-pooled.png}\hfill
  
\includegraphics[width=0.32\linewidth]{instrument-ratio.png}
  \caption{\figlead{The three views, on a $3{:}1$ pair of line families.}
  Left, the superposition as drawn. Middle, the pooled view: the slide
  average of Theorem~\ref{thm:sweep} along the schedule the station's
  recipe $(3,-1)$ asks for, integrated in closed form, so the bands are what
  a slow observer sees and nothing else. Right, the ratio view: the priced
  merit of Definition~\ref{def:merit} for the winning recipe at every pixel,
  dark where a beat must form. Rendered by the tool's own capture path,
  which the goldens of the regression suite also use.}
  \label{fig:instrument}
\end{figure}

Three of its views are this paper's theorems running
(Figure~\ref{fig:instrument}). The \emph{pooled} view is
Theorem~\ref{thm:sweep}: the slide average, integrated in closed form.
Along a slide every layer's coverage is its stroke profile, a window with
soft edges, evaluated at a distance that changes linearly with the slide,
so the coverage is a cubic polynomial on each stretch between corners and
the corners are known from the count; the average is therefore a sum of
integrals of polynomials with known endpoints, exact at every zoom, with no
sampling and no noise, and it runs at the frame rate (certified against a
$65{,}536$-point brute-force integral by \texttt{exactsweep.mjs}). The
\emph{ratio} view is Definition~\ref{def:merit}: the priced merit of the
winning recipe at every pixel, found by the lattice reduction of
Proposition~\ref{prop:selection}, drawn as a map that is dark where a beat
must form; for a first-order recipe it is the $\het$ of
\S\ref{sec:slow}. A \emph{square-law} toggle is
Theorem~\ref{thm:universality}: the drawing squared before it is averaged,
so that Corollary~\ref{cor:soft} can be watched. The strokes' soft edges
are a fixed width in pixels, so zooming out makes them a larger fraction
of the pitch, and the null of a $2{:}1$ pair at coarse duty one half, flat
under the linear pooled view at every zoom, reopens under the square-law
toggle as the zoom widens the edges.

The field language also runs the theory backwards. A picture with
brightness $g$ is a field: give a copy of a layer the index shift
$D=\tfrac12(1-g)$, read from the picture in the copy's own coordinates,
and the tent of Corollary~\ref{cor:cases} says what any pooling observer
sees of the pair. The strokes align where the picture is white and
interleave where it is black, so in register the overlay is the picture,
at half contrast, and nothing else (Figure~\ref{fig:inverse}). The amount
is a dial: a whole member of shift at black is invisible inside the figure,
so only its edges, where the shift passes through one half, stay dark, and
the same picture is drawn as an outline. Out of register the same two
layers are the classical band moir\'e with the picture riding on it. A
quarter of a pitch apart, the offset is one quarter outside the figure and
three quarters inside, which the tent cannot tell apart, and the picture
fades to its edges. Turned by under a degree, the bands of the rotation
carry it as a displacement, and as the turn closes the bands widen until
one of them is the whole frame and the picture is back. The construction
is one line because the count is the object: the picture is a slow
combination chosen by hand, and the base family is free.

It is also visual cryptography~\cite{NaorShamir1994Visual}, in the count's
currency, and the theory says what the shares must look like. A crisp
picture needs the difference of the two shifts to jump by one half at
every edge, so at least one layer jumps by a quarter there, and a smooth
layer cannot hide a coherent quarter-member jog along an edge: subtle,
crisp and hidden are two of three. What gives all three is a jitter that
is everywhere. Cut each layer into cells one pitch across the members and
about one pitch along them, and shift every cell's stroke by a quarter
member to one side or the other, at random, the same random signs on both
layers. On one of them, flip the sign of every cell inside the picture.
Each layer alone is then the same weave inside the picture and out, with
nothing to see, not even statistically; in register the jitters cancel
outside the picture and add to one half inside it, and the tent makes that
solid black at cell size (Figure~\ref{fig:inverse}, bottom row). A grey
cell flips with probability equal to its darkness, so grey levels come out
as dither. Desmedt and Van Le~\cite{DesmedtVanLe2000Moire} hid a secret in
the moir\'e of two random gratings by the same logic; here the shares are
counts, so the scheme works on any family the count reaches, and the
pooling theorem is the decoder. In \Moire\ the picture is the Picture
button of the field editor, and its Share row puts it on one layer, on two
layers as halves, or on two as a weave, the second layer being the
picture's unmodulated twin.

\begin{figure}[t]
  \centering
  \includegraphics[width=0.32\linewidth]{inverse-alone.png}\hfill
  
\includegraphics[width=0.32\linewidth]{inverse-aligned.png}\hfill
  \includegraphics[width=0.32\linewidth]{inverse-outline.png}\\[3pt]
  
\includegraphics[width=0.32\linewidth]{weave-keeper.png}\hfill
  
\includegraphics[width=0.32\linewidth]{weave-flipper.png}\hfill
  \includegraphics[width=0.32\linewidth]{weave-aligned.png}
  \caption{\figlead{The inverse problem, in the tool.} Top row, one layer
  carries the picture as the index shift $D=\tfrac12(1-g)$. Alone (left),
  the picture is in the lines. Over its unmodulated twin in register
  (middle), the strokes interleave where the picture is black and the
  overlay is the picture. With a whole member of shift at black and the
  edges softened (right), the inside is invisible and the picture is an
  outline. Bottom row, the weave: both layers jitter by a quarter member per
  cell with the same random signs, and one of them flips its cells inside
  the picture. The layer that keeps its signs (left) and the layer that
  flips them (middle) are, alone, the same weave everywhere; in register
  (right) the picture is solid at cell size. Rendered by the tool's capture
  path from scenes the field editor writes.}
  \label{fig:inverse}
\end{figure}

Every view is a measurement with a gate, in two different senses that
should not be confused. The count solver for the fold rung is checked
against a brute-force twin that tries every member, so its answers are
checked against truth. The renders are checked against golden images
through the real graphics device on every change, which catches any change
in what the tool draws but would preserve a wrong picture faithfully; the
goldens are trusted because the scripts of Table~\ref{tab:predictions}
check what they show against the theorems. And every number in this paper
is written into the text by the program that fails if an experiment does.
The tool, its source, and every script named in this paper are public at
\todo{URL, to be supplied by the author}; the systems paper describes the
engineering, and is where the proofs of \S\ref{sec:counting}'s closed
forms live.

\begin{tryit}{the three views}
In \Moire, make two layers of parallel lines at pitches $15$ and $5$ and
turn one by two degrees. The drawn view shows a fine texture with broad
bands in it; the pooled view keeps the bands and drops the texture; the
ratio view is dark exactly where the pooled view has bands. Now change the
coarse pitch to $16$: the bands change shape, the dark region moves before
they do, and the station table of \S\ref{sec:selection} says which recipe
you are looking at.
\end{tryit}
